The \textsc{AnchorBench} design supports evaluation of
inference-time mitigation strategies as a downstream use case.
While mitigation is not the primary contribution of this work,
we briefly describe six strategies implemented in the evaluation
framework, organized into prompt-level and multi-pass categories.

\paragraph{Prompt-level (single-pass).}
\textbf{B1} (Ignore-Numbers):
Prepends an instruction to ignore contextual numbers.
\textbf{B2} (Bias-Warning):
Prepends a warning about anchoring bias.
\textbf{B3} (Chain-of-Thought):
Requests step-by-step reasoning before answering.

\paragraph{Multi-pass.}
\textbf{SELFHELP} (Self-Debiasing):
The model rewrites the prompt to remove biasing numbers, then
answers the rewritten prompt.
\textbf{LATE\_BINDING} (Anchor-Late Reintroduction):
Pass~1 replaces the anchor with \texttt{[REDACTED]}; Pass~2
reveals the anchor and asks the model to update only if
justified.

\paragraph{Baseline.}
\textbf{B0} (No Mitigation): The unmodified anchored prompt.

\paragraph{Preliminary observations.}
In prior experiments on the v1 benchmark (which used only the
plausible relevance level), we observed that B2 (Bias-Warning)
provided the most consistent anchoring reduction across models,
while B3 (Chain-of-Thought) paradoxically amplified anchoring.
Full mitigation evaluation on the v2 benchmark with the
relevance axis is planned as future work and will measure
whether mitigations reduce irrelevant anchoring without
suppressing appropriate sensitivity to informative context.
